\section{Combinators}
\label{sec:regular-list-functions}

The elegant syntax of regular expressions is one of the main attractions of regular languages. Regular expressions have no variables or states, and they are compositional -- bigger regular expressions are built from smaller ones using a few simple combinators, namely union, concatenation and Kleene star. In this section, we describe an attempt at a similar syntax for regular functions. One could argue that the prime decompositions earlier in this book are such a syntax, see \cref{def:regular-functions}, with only one combinator, namely composition. However, those decompositions are motivated mainly by a mathematical desire to find a  minimal core of a given transducer class. The beauty of regular expressions is that they are both a minimal core, and are also  usable in everyday programming practice. The formalism in this section is intended to go in that direction, and to be more grounded in programming practice, by using  functions like 
\begin{align*}
 \myunderbrace{A^{**} \xrightarrow {\concatterm} A^*}{flatten a list of lists \\ into a list} \qquad \myunderbrace{\onetype + A \times A^* \xrightarrow {\consterm} A^*}{a list is constructed from \\ either an empty list, or \\ a head and a tail},
\end{align*}
which are familiar to any user of a functional programming language, such as  Haskell.

\paragraph*{Types.} The functions in the above examples are not based purely on strings, but have more structured types. For example, the domain of $\concatterm$ uses lists of lists, while the domain of $\consterm$ is constructed using products, co-products, lists and a unit type. Such types can be encoded into strings using separators, but we will prefer a more direct approach which does not use such representations. The types that we care about are defined in the following way. 
\begin{definition}
    [Types]\label{def:types} Types are expressions built using the following constructors:
\begin{align*}
    \myunderbrace{\onetype}{unit  type \\ which contains \\ a unique element $1 \in \onetype$} 
\qquad
\myunderbrace{A \times B}{product \\ type}
\qquad
\myunderbrace{A + B}{co-product \\ type}
\qquad
\myunderbrace{A^*}{list \\ type}.
\end{align*}
\end{definition}

Each type also represents a set in the natural way, and we will use the same notation for the type and the set that it represents. For example $\onetype^*$ represents list over the unit type; this set is in bijection with the natural numbers. Using multiple disjoint copies of the unit type we can represent finite sets, e.g.~$\onetype + \onetype$ represents a binary alphabet, and $\onetype + \onetype + \onetype$ represents a ternary alphabet. Formally speaking, for the ternary alphabet we should choose one of the two bracketings $\onetype + (\onetype + \onetype)$ or $(\onetype + \onetype) + \onetype$, but these two types are isomorphic and we will not  pay attention to the difference.

\paragraph*{Type-to-type functions.}
A type-to-type function is defined to be  any function 
\begin{align*}
A \xrightarrow f B
\end{align*}
where $A$ and $B$ are types. As we have explained above, types can be used to describe finite alphabets and lists over these alphabets, and hence string-to-string functions are a special case of type-to-type functions.  However, the opposite inclusion is also essentially true, since types can be represented using strings, and thus type-to-type functions can be represented as string-to-string functions.
As the alphabet for the string representations,  we use  an eight-letter alphabet that we denote by $8$ because it has eight letters: 
\begin{align*}
( \quad ) \quad [ \quad ] \quad , \quad \onelement \quad \leftterm \quad \rightterm.
\end{align*}
 The string representation for elements of types is defined by induction:  the unique element of $\onetype$ is represented by  $\onelement$; elements of $A + B$ are represented by $\leftterm\  a$ or $\rightterm\  b$; pairs are represented by $(a,b)$, and lists are represented by $[a_1,\ldots,a_n]$. Using this representation, we can  lift transducer classes to describe type-to-type functions.







\begin{definition}[Regular under string representation]\label{def:regular-functions-on-types-under-string-representation}
    A type-to-type function
    \begin{align*}
    A \xrightarrow f  B
    \end{align*}
    is called \emph{regular under string representation} if there is some regular string-to-string function $f'$ such that the following diagram commutes:
    \[
    \begin{tikzcd}
    A 
    \ar[r,"f"]
    \ar[d,"\text{string representation}"']
    &
    \ar[d,"\text{string representation}"]
    B
    \\
    8^* 
    \ar[r,"f'"']
    &
    8^*.
    \end{tikzcd}
    \]
     Similarly we define rational functions between types.
\end{definition}

\begin{myexample}[Parity under string representation]
    Consider the function 
    \begin{align*}
    \onetype^* \xrightarrow f \onetype + \onetype
    \end{align*}
    which returns the parity of the input list. We think of the codomain as representing the Booleans, with the first copy of $\onetype$ representing ``yes'' and the second copy representing ``no''. Here are two examples of input/output pairs for this function under string representation:
    \begin{align*}
    [\onelement,\onelement,\onelement] \mapsto \myunderbrace{\rightterm\onelement}{``no'' means odd} \qquad [\onelement,\onelement] \mapsto \myunderbrace{\leftterm\onelement}{``yes'' means even}.
    \end{align*}
    The function $f$ is regular under string representation. The transducer $f'$ counts the appearances of the letter $1$ modulo two, and outputs $\leftterm\onelement$ or $\rightterm\onelement$ accordingly. Note that not all input strings over the input alphabet $8$ are representations of the domain, but the transducer $f'$ will still produce some output  for such inputs. However, the input strings which are correctly formatted, i.e.~represent some element of the domain, form a regular language, as witnessed by the following automaton:
    \mypic{54}
Therefore, we could improve $f'$ so that it produces the empty string for inputs that are not correctly formatted.
\end{myexample}


\subsection{Regular terms}
\label{sec:regular-terms}
As we mentioned earlier in this section, when defining type-to-type functions, we would like to work directly on the types instead of using string representations as is done in \cref{def:regular-functions-on-types-under-string-representation}. This is the purpose of the following definition, in which type-to-type functions are built by applying four standard combinators to certain atomic functions. We begin by describing two atomic functions that are not completely standard. 

\begin{myexample}[Split]\label{ex:split}
    For types $A$ and $B$, the split function 
    \begin{align*}
    (A + B)^* \xrightarrow{\splitterm} A^* \times (B \times A^*)^*
    \end{align*}
    splits the input list according to appearances of elements from $B$, as illustrated in the following example where elements of $A$ are letters and elements of $B$ are red digits:
    \begin{align*}
    [a,b,c,\red 1,d,e,\red 2,\red 3,f,g,h] \mapsto ([a,b,c],[(\red 1,[d,e]),(\red 2,[]), (\red 3,[f,g,h])]).
    \end{align*}
    A defining property of the split function is that it is a bijection, and its inverse is the concatenation of all its entries in the left-to-right order. 
\end{myexample}

\begin{myexample}[Group prefix multiplication]
    For a finite group $G$, the group prefix multiplication function 
    \begin{align*}
    G^* \xrightarrow{\prefterm} G^*
    \end{align*}
    takes a list of elements from $G$ and returns the list of all prefix products. This means that the output list has the same length as the input list, and the $i$-th element of the output list is the result of applying the group operation  to the first $i$ elements in the input list. For example, if $G$ is $\set{0,1,2}$ with addition modulo 3, then we have
    \begin{align*}
    [0,1,2] \mapsto [0,1,0].
    \end{align*}
    Note that this function is parameterised not just by the underlying set of the group, but also its group operation. 
\end{myexample}

Apart from the above two examples, all other atomic functions are standard built-in functions of Haskell, and we believe that the combinators are also self-explanatory. Here is the definition.

\begin{definition}[Regular term] \label{def:regular-terms} A \emph{regular term} is any expresion that is built by applying the combinators from \cref{item:regular-list-functions-combinators} to the atomic terms in~\cref{item:regular-list-functions-fo-atomic}. In the following $A,B,C$ range over types, and $G$ ranges over groups whose underlying set is a finite type. 
    \begin{enumerate}
        \item \label{item:regular-list-functions-fo-atomic}     The atomic functions are:
    \begin{description}
        \item[Identity]  $A \xrightarrow {\idterm} A$.
        \item[Projections]  $A \times B \xrightarrow {\fstterm} A$ and $A \times B \xrightarrow {\sndterm} B$.
        \item[Co-projections] $A \xrightarrow {\leftterm} A + B$ and $B \xrightarrow {\rightterm} A + B$.
        \item[Diagonal and co-diagonal]  $A \xrightarrow {\diagterm} A \times A$ and   $A + A \xrightarrow {\codiagterm} A$.
        \item[Distributivity]  $A \times (B + C) \xrightarrow {\distrterm} (A \times B) + (A \times C)$. 
        \item[List constructor/deconstructor] $\onetype + A \times A^* \xrightarrow {\consterm} A^*$ and its inverse $\unconsterm$. 
        \item[Reverse]  $A^* \xrightarrow{\reverseterm} A^*$.
        \item[Concatenation] $A^{**} \xrightarrow{\concatterm} A^*$.
        \item[Split] $(A+B)^* \xrightarrow{\splitterm} A^* \times (B \times A^*)^*$.
        \item[Group prefix multiplication] $G^* \xrightarrow{\prefterm} G^*$.
    \end{description}
    
            \item \label{item:regular-list-functions-combinators} The combinators are composition of functions, as well as the following functoriality combinators that lift functions along the type constructors:
        \begin{align*}
        \frac{ A_1 \xrightarrow {f_1} B_1 \quad A_2 \xrightarrow {f_2} B_2}
        { A_1 \times A_2 \xrightarrow {f_1 \times f_2} B_1 \times B_2}
        \qquad 
                \frac{ A_1 \xrightarrow {f_1} B_1 \quad A_2 \xrightarrow {f_2} B_2}
        { A_1 + A_2 \xrightarrow {f_1 + f_2} B_1 + B_2}
        \qquad 
        \frac { A \xrightarrow f B}
        { A^* \xrightarrow {f^*} B^*}.
        \end{align*}
    \end{enumerate}
\end{definition}

Each term has syntax (an expression built from the atomic functions and combinators) and semantics (a type-to-type function), but we will not distinguish between the two, and we will use the same notation for both. 

\begin{theorem}\label{thm:regular-terms} A type-to-type function is regular under string representation if and only if it is defined by some regular term.
\end{theorem}

The rest of \cref{sec:regular-terms} is devoted to the proof of \cref{thm:regular-terms}.  We begin with the easy direction, which is
\begin{align*}
\text{defined by a regular term} 
\quad \Rightarrow \quad 
\text{regular under string representation.}
\end{align*}
This is a straightforward induction on the structure of the term. The induction basis says that all atomic terms are regular under string representation, which is easy to check, with the only  non-trivial part being  parsing the string representation of the input. Let us illustrate this on one example, which is distributivity.

\begin{lemma}
    For every types $A,B,C$, the distributivity function
    \begin{align*}
    A \times (B + C) \xrightarrow {\distrterm} (A \times B) + (A \times C)
    \end{align*}
    is regular under string representation.
\end{lemma}
\begin{proof}
    The input, if correctly formatted, has one of the forms $(a, \leftterm b)$ or $(a, \rightterm c)$, where $a$ is a string representation of some element in $A$, and similarly for $b$ and $c$.      The transducer that implements the distributivity function first checks if the input is correctly formatted, and then checks if it uses the left or right variant, and finally outputs the corresponding output. In all of these steps, we use the fact that parsing can be implemented by a finite automaton, as stated in the following claim, which is proved by a straightforward induction on the structure of the type. 
        \begin{claim}
        For every type $A$, the following is a regular language
        \begin{align*}
        \setbuild{ w \in 8^*}{$w$ represents an element of $A$}.
        \end{align*}
    \end{claim}
    Using the above claim, an automaton can check if the input is correctly formatted, and if so, it can check if the input uses the left or right variant. Also, the output can be produced, by identifying the separation between the two components of the input. 
\end{proof}

The remaining atomic functions are handled in a similar way,  and the combinators are also easily seen to preserve regularity, thus proving the induction step.

The rest of this proof is devoted to the more difficult implication 
\begin{align*}
\text{defined by a regular term} 
\quad \Leftarrow \quad 
\text{regular under string representation,}
\end{align*}
which says that the terms are expressively complete. Ultimately, we need to prove expressive completeness for all type-to-type functions, but we begin with the string-to-string case, which contains the essence of the proof. In the string-to-string case, we can profit from the prime decomposition theorems (which are, technically speaking, our formal definition of the transducer class). Since terms are syntactically closed under composition, it suffices to show that the following prime  string-to-string functions are definable by terms, which is done in Lemmas \ref{lem:terms-define-string-homomorphisms} - \ref{lem:terms-define-reversible} as described in the following table.
\begin{center}
    \begin{tabular}{ll}
    \cref{lem:terms-define-string-homomorphisms} & string-to-string homomorphisms \\
    \cref{lem:terms-define-append-hash} & appending an endmarker $w \mapsto w\#$ \\
        \cref{lem:terms-define-map-reverse-duplicate} & map reverse and map duplicate \\
    \cref{lem:terms-define-flip-flop} & flip-flop Mealy machines \\
    \cref{lem:terms-define-reversible} & reversible Mealy machines \\
    \text{symmetric} & right-to-left variants of prime Mealy machines
\end{tabular}
\end{center}
In the last line of the above table, the symmetry is resolved by the string reversal function, since a right-to-left Mealy machine can be decomposed as: (1) reverse; (2) a left-to-right Mealy machine; (3) reverse again. 
From the decomposition result on rational functions in \cref{thm:rational-primes} and the definition of regular functions, the above lemmas will imply that all regular string-to-string functions are definable by terms. We will then conclude the proof in \cref{lem:terms-define-string-representation} by extending the completeness from string-to-string functions to all type-to-type functions. 

Let us begin with the case of string-to-string homomorphisms, which is easy and also establishes some infrastructure that will be used in the other cases.



\begin{lemma}\label{lem:terms-define-string-homomorphisms}
        Every string-to-string homomorphism is definable by a regular term.
\end{lemma}
\begin{proof} We begin by deriving two useful combinators.

    \begin{claim}\label{claim:pairing-copairing}
        Functions definable by terms are closed under the following combinators, which are called \emph{pairing} and \emph{co-pairing}:
        \begin{align*}
        \frac{ A \xrightarrow {f_1} B_1 \quad A \xrightarrow {f_2} B_2}
        { A \xrightarrow {(f_1,f_2)} B_1 \times B_2}
        \qquad
        \frac{ A_1 \xrightarrow {f_1} B \quad A_2 \xrightarrow {f_2} B}
        { A_1 + A_2 \xrightarrow {\either {f_1} {f_2}} B}.
        \end{align*}
    \end{claim}
    \begin{proof}
        For pairing, we apply the diagonal followed by  $f_1 \times f_2$.  For co-pairing, we apply $f_1 + f_2$ followed by  co-diagonal.
    \end{proof}

    We now analyse functions with finite domains.
    We say that a type is of the form $\onetype + \ldots + \onetype$ if it is a co-product of finitely many copies of the unit type.  (As mentioned before, a formal description would need to specify the bracketing, but  all possible bracketings are isomorphic using regular terms.) The following claim shows that such types cover all finite types.
    \begin{claim}\label{claim:finite-type-bijection-disjoint-units}
        Every finite type admits a bijection, which is definable by a regular term, to a type of the form $\onetype + \ldots + \onetype$.
    \end{claim}
    \begin{proof}
        Induction on the structure of the finite type, with distributivity used in the case of products. It is worth noting that the proof needs two  kinds of distributivity:
\begin{align*}
A \times (B + C) \xrightarrow {\distrterm} (A \times B) + (A \times C)\\
(A + B) \times C \xrightarrow {\distlterm} (A \times C) + (B \times C).
\end{align*}
    Our atomic functions contain only the first kind, but the second kind can be derived from the first one, since we have a function for swapping coordinates:
\begin{align*}
A \times B \xrightarrow {\mathtt{swap}} B \times A,
\end{align*}
which can be derived using pairing and projections.
    \end{proof}
    \begin{claim}\label{claim:finite-domain-regular-list-function}
        Every type-to-type function with a finite domain  is definable by a regular term.
    \end{claim}
    \begin{proof}
        By \cref{claim:finite-type-bijection-disjoint-units}, we can assume that the domain is of the form $\onetype + \ldots + \onetype$. By using co-pairing, we can reduce to the case of a function with domain $\onetype$. We now use induction on the structure of the codomain. The interesting case is a function 
\begin{align*}
\onetype \xrightarrow f A^*
\end{align*}
with a list codomain. Suppose that the output is a list $[a_1,\ldots,a_n]$. A term defining the function is constructed as follows. If $n=0$, i.e.~the output is the empty list, then we can use the composition
\begin{align*}
\onetype \xrightarrow{\leftterm} \onetype + A \times A^* \xrightarrow {\consterm} A^*.
\end{align*}
If the output list is non-empty, then use pairing to construct the pair (head, tail) for the list, with both the head and tail constructed using the induction hypothesis. Finally, use the list constructor to construct the output list from the head and tail. (There is a nested induction here on the length of the output list.) 
\end{proof}

By combining the above two claims, we see that if $A$ and $B$ are finite types, then every function of type $A \to B^*$ is definable by a regular term. Every  string-to-string homomorphism  then arises by taking such a function $f$, and using the following composition
\begin{align*}
    A^* \xrightarrow{f^*} B^{**} \xrightarrow \concatterm B^*.
\end{align*}
\end{proof}

\begin{lemma}\label{lem:terms-define-append-hash}
    The function $w \mapsto w\#$ is definable by a term.
\end{lemma}
\begin{proof}
    When formalising this function, we think of $\#$ as being a fresh symbol. Therefore, we can think of its type as being 
    \begin{align*}
    A^* \to (A + \onetype)^*.
    \end{align*}
    (This type of thinking is justified by \cref{claim:finite-domain-regular-list-function}, which says that any function with a finite domain is definable by a term, which covers bijections between finite types.) By the symmetry principle, it will be easier to define a prepending function $w \mapsto \#w$. We begin by defining functions into the unit type (which are unique).

    \begin{claim}\label{claim:bang-definable}
         For every type $A$, the function  $A \xrightarrow ! \onetype$ is definable by a term. 
    \end{claim}
    \begin{proof}
        Induction on the structure of the type, with the identity used for the induction bases, projections used for products, co-pairing used for co-products, and the list deconstructor used for lists.
    \end{proof}
    Using the $!$ function from the above claim, we can define prepending as follows.
    \begin{align*}
    A^* \xrightarrow {((!;\rightterm),\leftterm^*)}   (A + \onetype) \times (A + \onetype)^* \xrightarrow {\rightterm;\consterm} (A + \onetype)^*.
    \end{align*}
\end{proof}

\begin{lemma}\label{lem:terms-define-map-reverse-duplicate}
   Map reverse and map duplicate are definable by terms.
\end{lemma}
\begin{proof}
    We begin by defining some helper functions:
            \begin{align*}
        \onetype & \xrightarrow \emptyterm A^*  & \leftterm; \consterm \\ 
        A & \xrightarrow \pureterm A^* & (\idterm, (!;\emptyterm));\rightterm;\consterm\\
        A^* \times A^* & \xrightarrow \binconcterm A^* & (\idterm \times \pureterm); \rightterm; \consterm; \concatterm .
        \end{align*}
    In the last line above,  we  use a derived derived combinator for applying two functions in parallel to separate coordinates of a product, which is defined by
    \begin{align*}
    f \times g = ((\fstterm;f),(\sndterm;g)).
    \end{align*}
    Finally, we define the inverse of $\splitterm$, which is a form of concatenation that was discussed earlier in Example~\ref{ex:split}. This inverse is the composition of 
    \begin{align*}
    A^* \times (B \times A^*)^* \xrightarrow{\leftterm \times (\rightterm \times \leftterm^*)^*} (A+B)^* \times ((A+B)^* \times (A+B)^*)^*
    \end{align*}
    followed by the below form of concatenation where $C = A+B$:
    \begin{align*}
    C^* \times (C \times C^*)^* \xrightarrow{ \idterm \times (\rightterm; \consterm)^*} C^* \times C^{**} \xrightarrow{\rightterm;\consterm} C^{**} \xrightarrow \concatterm C^*
    \end{align*}

    Equipped with the above helper functions, we can define map reverse and map duplicate, as required by the lemma.
    We assume that these functions are typed as 
    \begin{align*}
    (A + \onetype)^* \to (A + \onetype)^*.
    \end{align*}
    The term for map reverse is the following composition:
    \begin{align*}
    (A + \onetype)^* \xrightarrow{\splitterm} A^* \times (\onetype \times A^*)^* \xrightarrow{(\reverseterm \times (\idterm \times \reverseterm)^*)} A^* \times (\onetype \times A^*)^*,
    \end{align*}
    followed by the inverse of $\splitterm$. In the presence of these helper functions, map duplicate is defined in the same way as map reverse, except that we use $(\idterm,\idterm);\binconcterm$ instead of string reversal.
\end{proof}

\begin{lemma}\label{lem:terms-define-flip-flop}
        Every flip-flop Mealy machine is definable by a term.
\end{lemma}
\begin{proof}    
    Consider a flip-flop Mealy machine 
    \begin{align*}
    A^* \xrightarrow f B^*
    \end{align*}
    with states $Q$. We assume that the states are a finite type. 
    The input alphabet partitions into two parts: the identity letters which leave the state unchanged, and the reset letters which reset the state to a fixed state. To define this function using a term, we will split the input string according to appearances of the reset letters. The exact construction has several steps, which are described below.

    \begin{enumerate}
        \item Before splitting along the reset letters, we will need to annotate them with the states that they produce. This is implemented by a string-to-string homomorphism  
        \begin{align*}
        A^* \xrightarrow h (A+Q)^*
        \end{align*}
        in which the identity letters are left unchanged, and each  reset letter $a$ is replaced by the two-letter string $aq$ where $q$  is the state  after reading $a$.  
        \item In the second step, we apply $\splitterm$ to  the output of the first step in order to split along the states, thus yielding an element of
        \begin{align*}
         (w_0, [(q_1,w_1),\ldots,(q_n,w_n)]) \in  A^* \times (Q \times A^*)^*.
        \end{align*}
        Thanks to the way that the first step was defined, we know that: (1) each of the strings $w_0,\ldots,w_n$ contains at most one reset letter, and if it contains a reset letter then it is the last one; and  (2)  $q_i$ is the state of the Mealy machine after reading $w_0 \cdots w_{i-1}$. In particular, thanks to (1), we know that the state of the Mealy machine will not change while reading any of the strings $w_0,\ldots,w_n$ (except at the end, when the state no longer influences the output). We will leverage this in the following steps, by simulating the Mealy machine using a homomorphism. 
        \item Thanks to the way that the previous step was defined, the output produced by the Mealy machine on each pair $(q_i,w_i)$ is obtained by applying the function 
        \begin{align*}
                Q \times A^* & \xrightarrow g B^*\\ 
                (q,v) & \mapsto h_q(v)
        \end{align*}
        where $h_q$ is the string-to-string homomorphism that replaces each letter by the corresponding output of the Mealy machine assuming that the current state is $q$.  Using \cref{lem:terms-define-string-homomorphisms}, as well as distributivity and co-pairing, we can show that $g$ is definable by a term.. (The variant of distributivity is the left-hand one that was given in the proof of \cref{claim:finite-type-bijection-disjoint-units}.)  For the initial string $w_0$, we can simply apply the homomorphism $h_{q_0}$, where $q_0$ is the initial state of the Mealy machine.
        \item As a result of the previous step, we get an element of 
            \begin{align*}
            B^* \times (B^*)^*
            \end{align*}
            which should be concatenated to get the output string. This concatenation can be done by a term. 
    \end{enumerate}
\end{proof}

\begin{lemma}
    \label{lem:terms-define-reversible}
    Every reversible Mealy machine is definable by a term.
\end{lemma}
\begin{proof}
        Consider a reversible  Mealy machine 
    \begin{align*}
    A^* \xrightarrow f B^*
    \end{align*}
    with states $Q$.  Let us also choose some arbitrary group structure on the input alphabet $A$, say a cyclic group. 

    \begin{itemize}
        \item  Let $G$ be the group of permutations of the finite set $Q$; this is a finite group and we can describe its underlying set using some type. In the first step we apply the letter-to-letter homomorphism
        \begin{align*}
        A^* \to (A \times G)^*
        \end{align*}
        where each letter is paired with its state transformation.
        \item Let the output of the first step be
        \begin{align*}
        [(a_1,g_1), \ldots, (a_n,g_n)] \in (A \times G)^*. 
        \end{align*}
        Thanks to the group structure on $A$,  the alphabet $A \times G$ can be seen as a group. In the second step, we apply group prefix multiplication.
        \item The result of the previous step is 
        \begin{align*}
        [(a'_1,g'_1),  \ldots, (a'_n, g'_n)] \in (A \times G)^*,
        \end{align*}
        where $a'_i$ is the product of the first  $i$  input letters (according to our arbitrary group structure) and $g'_i$ is the state transformation induced by the first $i$ input letters. Recall the  delay machine from Example \ref{ex:delay}, which is a flip-flop Mealy machine (and hence definable by a term thanks to \cref{lem:terms-define-flip-flop}). Using this machine, we can label each pair with the previous one, so that each pair in the list becomes a four-tuple $(a'_i,g'_i,a'_{i-1},g'_{i-1})$, with the edge cases of  $a'_0$ and $g'_0$ being defined to be  the identity elements of the respective groups. 
        \item Using group inverses we can recover the original input letters, because 
        \begin{align*}
        a_i = (a'_{i-1})^{-1} \cdot a'_i.
        \end{align*}
        Therefore, for each item in the previous list, we can compute the pair $(a_i,g'_{i-1})$ which describes the original input letter, and the state transformation before reading it. This is all the information needed to get the corresponding output letter in the Mealy machine, and hence a string-to-string homomorphism can be used to get the output string.
    \end{itemize}
\end{proof}

We have thus proved that all prime regular string-to-string functions are definable by terms, and hence all regular string-to-string functions are definable by terms. It remains to lift the result to type-to-type functions. Consider  a type-to-type function 
\begin{align*}
A \xrightarrow f B
\end{align*}
that is regular under string representation. To define this function by a regular term, we compose the terms for: (a) the string representation of $A$; (b) the regular string-to-string function that defines $f$ on string representations; and (c) the inverse of the string representation of $B$.  Step (b) is already done, and we now need to show that steps (a) and (c) can be defined by terms. This will be proved in the following lemma.  The inverses in the lemma are understood as one-sided inverses, i.e.~a function such that the composition 
\begin{align*}
A \xrightarrow{\text{string representation}}
8^*
\xrightarrow{\text{inverse of string representation}}
A
\end{align*}
is the identity on $A$.


\begin{lemma}\label{lem:terms-define-string-representation}
    For every type, its string representation and a one-sided inverse  can be defined by terms.
\end{lemma}
\begin{proof}
    Induction on the structure of the type. We only do the induction step for a product type $A \times B$, with the other cases being similar. The string representation is easy, and we focus on its inverse. If correctly formatted, the input string is of the form $(a,b)$, where $a$ and $b$ are string representations of elements of $A$ and $B$, respectively. Consider the rational function 
    \begin{align*}
    8^* \xrightarrow f (8 + \onetype)^*
    \end{align*}
    which identifies the separating comma between $a$ and $b$ (there might be other commas in the representations of $a$ and $b$), and writes it using a fresh symbol. This function is defined by a term, since we have already proved expressive completeness for terms in the case of string-to-string functions. Apply the function $f$, followed by split, leading to an element of the type
\begin{align}\label{eq:split-product-unrepresentation}
    8^* \times (\onetype \times 8^*)^*.
\end{align}
We know that if the original input was correctly formatted, then the list on the second coordinate will have exactly one element. We will now extract it using a version of the head operation.
\begin{claim}
    \label{claim:head}
    For every type $A$ there is a term  $A^* \xrightarrow \headterm A$ which returns the first element when the input list is nonempty (and an arbitrary element when it is empty). 
\end{claim}
\begin{proof}
    The difficulty is inventing a value when the input list is empty. By a simple induction on the structure of the type $A$, we show that there is a term
    \begin{align*}
    \onetype \xrightarrow{\inventterm} A.
    \end{align*}
    Using this term, the claim is proved as follows
    \begin{align*}
   \headterm =  \unconsterm; \either  \inventterm \fstterm
    \end{align*}
\end{proof}
By applying $\headterm$  to the second coordinate in~\eqref{eq:split-product-unrepresentation}, we get an element of 
\begin{align*}
8^* \times (\onetype \times 8^*).
\end{align*}
By projecting away the $\onetype$, removing the outer brackets (which are the first letter of the first $8^*$ and the last letter of the second $8^*$), and applying the induction assumption to both coordinates, we get the desired element of $A \times B$. This completes the induction step for the product type, and the other cases are similar.
\end{proof}


\subsection{Rational terms}
\label{sec:rational-terms}
In \cref{def:regular-functions-on-types-under-string-representation}, we defined both rationality and regularity under string representations, but we characterised only the regular functions. We now isolate the fragment that defines the rational functions.  Clearly we need to remove the $\reverseterm$ function, since string reversal is not rational as we have shown in Example \ref{ex:string-reversal-not-rational}. We also need to remove the diagonal since it enables string duplication, which is not rational as we have shown in \cref{exer:duplication-not-rational}. However, once we remove string reversal and  the diagonal (and thus pairing), we need to add some  some  atomic functions and combinators that were previously definable using pairing or the symmetry accorded by list reversal. The following theorem gives the exact characterisation. 



\begin{theorem}\label{thm:rational-terms}
    A type-to-type function is rational under string representation if and only if it is defined by the following variant of regular terms: 
    \begin{enumerate}
    \item  We keep the same atomic functions and combinators as in \cref{def:regular-terms}, except that we remove  $\reverseterm$ and $\diagterm$  and we add the following:
        \begin{description}
        \item[Adjoining units]  $A  \xrightarrow {\appendoneterm} A \times \onetype$ and $A  \xrightarrow {\prependoneterm} \onetype \times A$.
        \item[Left distributivity]  $(A + B) \times C \xrightarrow {\distlterm} (A \times C) + (B \times C)$. 
        \item[List deconstructor on the right] $A^* \xrightarrow {\unsnocterm} \onetype + A^* \times A$.
    \end{description}
    
            \item We keep the same combinators as in \cref{def:regular-terms}.
    \end{enumerate}
\end{theorem}
\begin{proof}
    We simply follow the proof of \cref{thm:regular-terms}, while paying attention to the atomic functions that are used.
    We use the name \emph{rational terms} for the variant of regular terms described in the theorem.   The proof that all rational terms are rational under string representation is straightforward, and follows the same kind of inductive structure as in the regular case.  For the converse implication, the necessary changes to the proof of \cref{thm:regular-terms} are listed below:
    \begin{itemize}
    \item \cref{lem:terms-define-string-homomorphisms}  about string-to-string homomorphisms. In \cref{claim:finite-type-bijection-disjoint-units}, w need to use both kinds of distributivity, but these are now included as atomic functions. Also, after applying distributivity, each type admits a bijection with a disjoint union of products of the form $\onetype \times \cdots \times \onetype$. Each such product admits a bijection with $\onetype$, however the direction $\onetype \to \onetype \times \cdots \times \onetype$ needs to adjoin units (previously it could have been implemented using the diagonal). Also, pairing was used in \cref{claim:finite-domain-regular-list-function}, but again all that was needed was adjoining units.
    \item \cref{lem:terms-define-append-hash} about appending an endmarker $w \mapsto w\#$. In the proof, we use symmetry (which was available thanks to string reversal) to define a prepending function $w \mapsto \#w$. In the rational case, this is no longer available. However, we can use binary concatenation to do this. Recall the implementation of binary concatenation at the beginning of the proof of \cref{lem:terms-define-map-reverse-duplicate}. That implementation used pairing to implment $\pureterm$, but this can be worked around by adjoining units, implementing $\pureterm$ as
    \begin{align*}
    A 
    \xrightarrow {\appendoneterm}
    A \times \onetype 
    \xrightarrow {\idterm \times \emptyterm}
    A \times A^* 
    \xrightarrow {\rightterm; \consterm}
    A^*
    \end{align*}
    , but only to adjoin an empty list, and the same effect can obtained by adjoining a unit followed by the combinator $\times$. This way 
    \item    \cref{lem:terms-define-map-reverse-duplicate} about map reverse and map duplicate. This lemma is no longer needed, since these are not rational. 
    \item  \cref{lem:terms-define-flip-flop,lem:terms-define-reversible} about flip-flop and reversible Mealy machines did not use string reversal or the diagonal. 
    \item The right-to-left variants of the prime Mealy machines were previously implemented using string reversal and the left-to-right variants. String reversal is no longer available in the rational case. For reversible Mealy machines, there is nothing to do, since we have shown in \cref{exer:no-reverse-reversible} that right-to-left reversible Mealy machines are not needed in the presence of the other machines. Therefore, we are left with the right-to-left variant of flip-flop Mealy machines.  We can simply do the same proof again in the right-to-left variant. To do this, we need symmetric versions of the atom functions, namely distributivity, the list constructor/deconstructor, and the group prefix multiplication. These are all included in the list of atomic rational functions, except the list constructor, which can be implemented using binary concatenation. 
\end{itemize}

\end{proof}


\exerciseheader

\exer{\label{exer:if-then-else} Recall that we think of the type $\onetype + \onetype$ as representing Booleans. Show that functions definable by regular (or rational) terms are closed under the following conditional construction
\begin{align*}
a \in A \mapsto 
\begin{cases}
    f_1(a) & \text{if $f(a)$ is true} \\
    f_2(a) & \text{if $f(a)$ is false}
\end{cases}
\qquad \text{ where }
A \xrightarrow f \onetype + \onetype
\text{ and }
A \xrightarrow {f_1,f_2} B.
\end{align*}
} { For regular terms, we use the following composition:
\begin{align*}
A \xrightarrow{(\idterm, f)} A \times (\onetype + \onetype)  \xrightarrow{\distrterm} (A \times \onetype) + (A \times \onetype) \xrightarrow {\fstterm \times \fstterm} A + A \xrightarrow {\either {f_1} {f_2}} B.
\end{align*}
For rational terms we cannot use the pairing in the first step. However, we can an indirect solution: by \cref{thm:rational-terms}, the functions defined by rational terms are exactly those that are rational under string representaiton, and the latter class of functions is closed under the conditional construction.
}


\exer{\label{exer:identity-function-replacement}Show that the general version $A \xrightarrow \idterm A$ of the identity function can be replaced by its special case of $A = \onetype$, and the expressive power of the regular terms is not affected. In fact, even the special case is not needed. }{
 By induction on the structure of the type $A$, we can lift the identity function from $\onetype$ to arbitrary types. We can implement the identity function on $\onetype$ as follows:
\begin{align*}
\onetype \xrightarrow{\rightterm} \onetype + \onetype  \xrightarrow{\codiagterm} \onetype.
\end{align*}
}



\exer{\label{exer:group-prefix-multiplication-isomorphism}Show that group prefix multiplication is an isomorphism.}
{ 
    Suppose that the output is 
    \begin{align*}
    [g_1,\ldots,g_n].
    \end{align*}
    Then the $i$-th element of the input is $g_{i-1}^{-1} g_i$, with the convention that $g_0$ is the identity element of the group.  The inverse is definable by a regular term. A curious question: is whether every bijective regular function has a regular inverse. 
}

\exer{\label{exer:windows-of-size-two} Write a regular term of type $A^* \to (A \times A)^*$ which outputs all windows of size two in the input list, as in the following example 
\begin{align*}
[1,2,3,4] \mapsto [(1,2),(2,3),(3,4)].
\end{align*}
}{ 
    Consider a function 
    \begin{align*}
    A &\xrightarrow f (A+A)^* \\
    a & \mapsto [\leftterm a, \rightterm a].
    \end{align*}
If we apply $f^*$ followed by $\concatterm$, then the output looks like this 
\begin{align*}
[\leftterm 1, \rightterm 1, \leftterm 2, \rightterm 2, \leftterm 3, \rightterm 3, \leftterm 4, \rightterm 4].
\end{align*}
We now apply $\splitterm$ to this output, which yields a value like this
\begin{align*}
([\leftterm 1],[(\rightterm 1, [\leftterm 2]), (\rightterm 2, [\leftterm 3]), (\rightterm 3, [\leftterm 4]), (\rightterm 4, [])]).
\end{align*}
Now project to the second coordinate, giving a value like this: 
\begin{align*}
[(\rightterm 1, [\leftterm 2]), (\rightterm 2, [\leftterm 3]), (\rightterm 3, [\leftterm 4]), (\rightterm 4, [])].
\end{align*}
Remove the last list element, and apply $(\idterm, \headterm)^*$, leading to the correct output.
}





\knowledgeconfigure{quotation=false}% tikzcd "..." labels need catcode-12 quotes inside \exer

\exer{\label{exer:natural-identity} Consider a family of functions 
\begin{align*}
A \xrightarrow {f_A} A \qquad \text{for each type $A$},
\end{align*}
which is natural in the sense that the diagram
\[
\begin{tikzcd}[ampersand replacement=\&]
A \arrow[r, "f_A"] \arrow[d, "g"'] \& A \arrow[d, "g"] \\
B \arrow[r, "f_B"'] \& B
\end{tikzcd}
\]
commutes for every function $g: A \to B$. Show that each $f_A$ must be the identity function on $A$.
}
{
    We do not even need to compare the functions for different types $A$. Consider some type $A$  and a function $f : A\to A$ 
    which is natural in the sense that $f;g = g;f$ holds for every function $g : A \to A$. We will show that $f$ must be the identity function. Indeed, let $a \in A$ and take a function $g$ which maps all elements to $a$. Then we have 
    \begin{align*}
       a \quad  \stackrel{\text{definition of $g$}} = \quad    
       g(f(a))
       \quad  \stackrel{\text{naturality}} = \quad 
        f(g(a)) 
      \quad  \stackrel{\text{definition of $g$}} = \quad   
        f(a) .
    \end{align*}
}

Let us continue the discussion of naturality. For the purposes of these exercises, define a \emph{functor} to be an expression  that is built from a type variable $A$ by using the type constructors from \cref{def:types}, such as the following expressions \begin{align*}
F(A) = A^* \qquad \text{or} \qquad   G(A) =  \onetype + A \times A^*.
\end{align*}
A functor can be applied to both types and to functions. For example, if we apply the functor $G$ above to some function $f : A \to B$, then we get a function
\begin{align*}
    \onetype + A \times A^* \xrightarrow {G(f)} \onetype + B \times B^*.
\end{align*}
Suppose that we have two functors $F$ and $G$ and a family of functions 
\begin{align*}
F(A) \xrightarrow {f_A} G(A) \qquad \text{for every type $A$}.
\end{align*}
We say that the family is a \emph{natural transformation} if for every function $g : A \to B$ the following diagram commutes:
\[
\begin{tikzcd}[ampersand replacement=\&]
F(A) \arrow[r, "f_A"] \arrow[d, "G(g)"'] \& G(A) \arrow[d, "G(g)"] \\
F(B) \arrow[r, "f_B"'] \& G(B)
\end{tikzcd}
\]



\exer{ \label{exer:head-operation-natural} Show that there is no natural transformation of type $A^* \to A$. In particular, the   $\headterm$ operation from \cref{claim:head} is not  a natural transformation.
}{
    Take any type $A$ and apply $f_A$ to the empty list, yielding some $a \in A$.  Take now any type $B$ that has at least two elements, e.g.~$B = \onetype + \onetype$, 
    and consider two functions 
    \begin{align*}
    A \xrightarrow {g_1, g_2} B
    \end{align*}
    which have different outputs when applied to $a$. (Such a functions can always be found, and even they can be defined by regular or even rational terms.) From naturality, we would need to have 
    \begin{align*}
    g_1(f_A([])) = f_B((g_1)^*([])) = f_B([]) = f_B((g_2)^*([])) = g_2(f_A([])),
    \end{align*}
     a contradiction.
}


\exer{\label{exer:natural-list-function} Show that there are uncountably many natural transformations of type $A^* \to A^*$.
}{
    Choose some function 
    \begin{align*}
     \set{1,2,\ldots} \xrightarrow \alpha  \set{1,2,\ldots}^*
     \quad \text{such that $\alpha(n)$ uses only letters from $\set{1,\ldots,n}$},
    \end{align*}
    which can be done in uncountably many ways. To this function we will associate a natural transformation, which is different for each choice of $\alpha$, thus proving part (b). The idea is that $\alpha$ tells us how the output list is constructed from the input list, by telling which input list element will be used for which output list element. Having fixed $\alpha$, define a family of functions as follows. The empty list is mapped to itself, and nonempty lists are mapped as follows:
    \begin{align*}
   [a_1,\ldots,a_n] \quad \stackrel{f_A} \mapsto \quad  [a_{\alpha(1)},\ldots,a_{\alpha(n)}].
    \end{align*}
    One can check that this family is natural, and that different choices of $\alpha$ lead to different families, thus proving part (b). One can, in fact, shown that all natural transformations arise in this way. Indeed, take some natural transformation. To define $\alpha$, we check what this natural transformation does for some infinite type $A$. Take such a type, and consider the output of $f_A$ on a non-repeating list $[a_1,\ldots,a_n]$. By naturality, the output list can only contain elements from the input list, and the pattern in which these elements appear can be used to define $\alpha(n)$. 
}

\exer{\label{exer:natural-transformations-from-A}  Show that all natural transformations of type $A \to A^*$ are regular (i.e.~regular for each type). }{
    By similar reasoning as in the previous exercise, one can show that such a natrual transformation is uniquely determined by the length of the output list, and it must simply repeat the input element that many times. 
}



\exer{\label{exer:natural-transformations-from-regular-terms} The notion of natural transformation is easily extended to functors that use multiple type variables. With the exception of group prefix multiplicaiton, the atomic functions in \cref{def:regular-terms}  are indexed by types, and hence they can potentially be natural transformations. Show that they are indeed natural transformations. 
}{
 A routine inspection.
}

\exer{\label{exer:natural-transformations-list-free} Consider a natural transformation, possibly with more than one type parameter, where the domain functor does not use the list constructor. Show that such a natural transformation is definable by a regular term. \footnote{As a corollary, we could more succinctly say that the atomic functions in \cref{def:regular-terms} are all natural transformations, as well as the functions that have lists in their domain, which are group prefix multiplication, $\reverseterm, \consterm, \unconsterm, \concatterm$ and $\splitterm$.}
}{
    Suppose first that the domain type is a product 
    \begin{align*}
    F(A_1,\ldots,A_n) = A_{i_1} \times \cdots \times A_{i_k},
    \end{align*}
    in which the types might be repeated.  By a similar analysis to to the one in \cref{exer:natural-list-function}, we can show that the natural transformation constructs the output in some fixed shape, with pointers to the input types (if a type variable appears multiple times in the input type, then the pointer must choose which one to use). For example, if we consider a natural transformation of type
    \begin{align*}
    A \times A \times B \to ((A+B)^* \times A)^* + A,
    \end{align*}
    then the outer $+$ must always be resolved in the same way, every output element from $B$ must be replaced by the input element from $B$, and every output element from $A$ must be replaced by either the first or the second input element from $A$, with the choice not depending on the type or input value.

    Let us now consider the general case. One can easily show that a natural transformation are closed under composition. Since distributivity is a natural transformation it follows that if we have  a natural transformation with domain 
    \begin{align*}
    A \times (B + C),
    \end{align*}
    then precomposing it with distributivity must yield a natural transformation with domain
    \begin{align*}
    A \times B + A \times C.
    \end{align*}
    As a result, we can use distributivity to reduce the domain of a natural transformation to a disjoint union of products of types. If the input type is a disjoint union, then the natural transformation is necessarily obtained by combining two natural transformations, one for each summand. Therefore, we can reduce to the case of a natural transformation whose domain is a product of type variables, which we have already analysed.
}


\knowledgeconfigure{quotation=true}